\documentclass[a4paper,11pt]{article}

\usepackage[utf8]{inputenc}
\usepackage[T1]{fontenc}
\usepackage[slovak]{babel}
\usepackage{amsmath,amssymb,amsfonts}
\usepackage{bm}
\usepackage{geometry}
\usepackage{enumitem}
\usepackage{listings}
\usepackage{xcolor}
\usepackage{hyperref}
\usepackage{booktabs}
\usepackage{parskip}

\geometry{margin=2.5cm}

\hypersetup{
  colorlinks=true,
  linkcolor=blue!60!black,
  citecolor=blue!60!black,
  urlcolor=blue!60!black
}

\lstset{
  language=Python,
  basicstyle=\ttfamily\small,
  keywordstyle=\color{blue!70!black},
  commentstyle=\color{green!50!black},
  stringstyle=\color{red!60!black},
  backgroundcolor=\color{gray!5},
  frame=single,
  framerule=0.4pt,
  rulecolor=\color{gray!40},
  breaklines=true,
  showstringspaces=false,
  columns=flexible,
  aboveskip=8pt,
  belowskip=8pt
}

\title{\textbf{Technická správa: EM algoritmus pre rekonštrukciu obrazu\\pomocou 2D Gaussovských zmesí}}
\author{Martin Hornáček}
\date{}

\begin{document}
\maketitle

\noindent
Tento dokument podrobne popisuje implementáciu rekonštrukcie obrazu pomocou
Expectation-Maximization (EM) algoritmu s~Gaussian Mixture Modelom (GMM).
Referenčný kód sa nachádza v~súboroch \texttt{run\_em\_algorithm.py} a~\texttt{utils/em\_utils.py}.

% ==========================================================================
\section{Načítanie a predspracovanie vstupného obrázka}
\label{sec:loading}
% ==========================================================================

Obrázok sa načítava pomocou knižnice PIL (\texttt{Image.open}), následne sa
konvertuje na RGB a~zmenší na rozmer špecifikovaný v~konfigurácii (typicky
$256\times256$ pixelov).

\begin{lstlisting}
# Nacitanie a konverzia
img = Image.open(img_path).convert('RGB')

# Zmena velkosti (riadena config suborom)
img = img.resize(tuple(config['dataset']['image_size']), Image.LANCZOS)

# Konverzia na numpy pole (float32, normalizovane do [0, 1])
gt_np = np.array(img).astype(np.float32) / 255.0
\end{lstlisting}

\noindent
Referencia v~kóde: \texttt{run\_em\_algorithm.py:218--220}.

% ==========================================================================
\section{Princíp GMM a~EM algoritmu}
\label{sec:math}
% ==========================================================================

Táto sekcia stručne vysvetľuje princípy Gaussian Mixture Modelu a~EM
algoritmu.  Presné matematické formulácie sú uvedené v~sprievodnom
článku (Section~2.1).

% --------------------------------------------------------------------------
\subsection{Gaussian Mixture Model --- intuitívny prehľad}
\label{subsec:gmm_def}
% --------------------------------------------------------------------------

Gaussian Mixture Model (GMM) modeluje obraz ako zmes $K$ Gaussiánov.
Každý Gaussián si možno predstaviť ako rozmazanú farebnú škvrnu (\emph{splat})
umiestnenú niekde v~obraze.  Každý komponent je definovaný tromi
parametrami:
\begin{itemize}[nosep]
  \item \textbf{Stred} $\bm{\mu}_k$ --- kde je škvrna umiestnená
        (pozícia + farba),
  \item \textbf{Kovariančná matica} $\bm{\Sigma}_k$ --- tvar a~orientácia
        škvrny (kruhová, eliptická, natočená),
  \item \textbf{Váha} $\pi_k$ --- relatívna dôležitosť škvrny
        (koľko pixelov modeluje).
\end{itemize}

V~našej implementácii sú vstupné dáta päťrozmerné: pozícia pixelu
$(x, y)$ a~farba $(r, g, b)$, kde všetky hodnoty sú normalizované
do~$[0,\,1]$.  Každý pixel je teda bod v~5D priestore a~GMM
sa snaží nájsť $K$ Gaussiánov, ktoré čo najlepšie pokryjú rozloženie
týchto bodov.

\begin{lstlisting}
# Kazdy pixel sa stane 5D bodom: [x_norm, y_norm, r, g, b]
def _prepare_data(img_np):
    h, w, _ = img_np.shape
    yy, xx = np.meshgrid(np.arange(h), np.arange(w), indexing='ij')
    coords = np.stack([xx.ravel(), yy.ravel()], axis=1)
    pixels = img_np.reshape(-1, 3)
    coords_norm = coords / np.array([[w, h]], dtype=np.float32)
    data = np.concatenate([coords_norm, pixels], axis=1)
    return data, h, w  # data ma tvar (H*W, 5)
\end{lstlisting}

% --------------------------------------------------------------------------
\subsection{EM algoritmus --- ako to funguje}
\label{subsec:em_alg}
% --------------------------------------------------------------------------

EM (Expectation-Maximization) algoritmus iteratívne hľadá najlepšie
parametre pre Gaussiány.  Funguje v~dvoch striedajúcich sa krokoch:

\begin{enumerate}[nosep]
  \item \textbf{E-krok (priradenie):} Pre~každý pixel sa vypočíta,
        s~akou pravdepodobnosťou patrí ku~každému Gaussiánu.
        Pixel blízko stredu veľkého Gaussiánu dostane vysokú zodpovednosť
        voči tomuto komponentu, zatiaľ čo vzdialené pixely dostanú takmer
        nulovú.
  \item \textbf{M-krok (aktualizácia):} Na základe priradených pixelov
        sa prepočítajú parametre každého Gaussiánu --- stred sa posunie
        do ťažiska priradených pixelov, kovariancia sa prispôsobí ich
        rozptylu a~váha sa upraví podľa počtu priradených pixelov.
\end{enumerate}

Tieto dva kroky sa opakujú, kým sa model nestabilizuje (konvergencia)
alebo sa nedosiahne maximálny počet iterácií.  Dôležitou vlastnosťou
EM algoritmu je, že kvalita modelu sa v~každej iterácii zaručene
nezhoršuje.

V~implementácii používame \texttt{sklearn.mixture.GaussianMixture}, čo
znamená, že samotné EM výpočty sú zapuzdrené v~knižnici:

\begin{lstlisting}
# Fitovanie GMM pomocou scikit-learn
gmm = GaussianMixture(
    n_components=K,          # pocet Gaussianov (napr. 256)
    covariance_type='full',  # plne kovariancne matice
    max_iter=500,            # max. pocet EM iteracii
    init_params='kmeans',    # inicializacia pomocou K-means
    random_state=42
)
gmm.fit(data_5d)  # data_5d ma tvar (H*W, 5)

# Vysledne parametre:
# gmm.means_       : (K, 5)    -- stredy [x, y, r, g, b]
# gmm.covariances_ : (K, 5, 5) -- kovariancne matice
# gmm.weights_     : (K,)      -- vahy (sucet = 1)
\end{lstlisting}

% --------------------------------------------------------------------------
\subsection{Rovnica renderovania}
\label{subsec:render_eq}
% --------------------------------------------------------------------------

Po naučení GMM sa obraz rekonštruuje renderovaním každého Gaussiánu ako
2D farebnej škvrny.  Pre pozíciu pixelu $\mathbf{p} = (x, y)$:
\begin{equation}
  I(\mathbf{p}) = \sum_{k=1}^{K} \pi_k \cdot \alpha_k(\mathbf{p})
                  \cdot \mathbf{c}_k,
  \label{eq:render}
\end{equation}

\noindent
Táto rovnica hovorí, že farba každého pixelu je súčtom príspevkov
všetkých $K$ Gaussiánov.  Jednotlivé členy majú nasledovný význam:

\begin{itemize}
  \item $\pi_k$ --- \textbf{váha Gaussiánu}: určuje, ako dôležitý je
        daný Gaussián v~celkovej zmesi.  Gaussián s~vyššou váhou
        pokrýva väčšiu oblasť obrazu a~jeho farebný príspevok je
        silnejší.

  \item $\alpha_k(\mathbf{p})$ --- \textbf{priestorový profil}:
        hovorí, ako blízko je pixel $\mathbf{p}$ k~stredu Gaussiánu.
        V~strede má hodnotu~1, smerom od stredu klesá podľa tvaru
        kovariančnej matice (kruhovo alebo elipticky).
        Technicky ide o~PDF Gaussiánu normalizovanú na maximálnu
        hodnotu~1 (\emph{peak-normalizácia}).

  \item $\mathbf{c}_k$ --- \textbf{farba Gaussiánu}: RGB hodnota
        naučená EM algoritmom.  Zodpovedá priemernej farbe pixelov
        priradených danému Gaussiánu (extrahovaná zo stredov:
        \texttt{gmm.means\_[k, 2:5]}).
\end{itemize}

\noindent
Výsledná farba pixelu je teda vážený súčet farebných príspevkov
všetkých Gaussiánov.  Gaussiány blízko pixelu prispievajú viac
(vďaka $\alpha_k$), dôležitejšie Gaussiány prispievajú viac
(vďaka $\pi_k$) a~každý prináša svoju farbu $\mathbf{c}_k$.
Finálna hodnota sa oreže do~$[0, 1]$.

\textbf{Analógia:} Každý Gaussián je ako sprej s~určitou farbou.
Váha $\pi_k$ udáva intenzitu striekania, $\alpha_k$ udáva, kam
dopadá farba (hustejšie v~strede, slabšie na okrajoch),
a~$\mathbf{c}_k$ je farba spreja.  Výsledný obraz je plátno,
na ktoré bolo nastriekaných $K$ takýchto farebných škvŕn.

Tento prístup zodpovedá $\alpha$-blendingu v~3D Gaussian Splatting,
ale bez triedenia podľa hĺbky --- v~2D všetky Gaussiány prispievajú
aditívne.

% ==========================================================================
\section{Príprava dát a~EM aproximácia obrazu}
\label{sec:fitting}
% ==========================================================================

Po načítaní obrázka do numpy poľa sa každý pixel transformuje do 5D
reprezentácie $[x, y, r, g, b]$, kde $x$ a~$y$ sú normované súradnice
pixelu a~$r, g, b$ sú normalizované hodnoty farieb.  Táto 5D matica
(\texttt{data\_5d}) je vstupom pre EM algoritmus.

\subsection{Príprava dát}

Funkcia \texttt{\_prepare\_data()} transformuje každý pixel na~5D vektor
$[x_{\mathrm{norm}}, y_{\mathrm{norm}}, r, g, b]$.  Súradnice sú normalizované
delením šírkou a~výškou obrazu, takže všetky hodnoty ležia v~$[0,1]$.
Výsledné pole má tvar $(H \cdot W,\; 5)$.

\begin{lstlisting}
# Priklad: obraz 256x256 = 65536 pixelov
data_5d, h, w = _prepare_data(gt_np)
print(data_5d.shape)  # (65536, 5)
print(data_5d[0])     # napr. [0.0, 0.0, 0.54, 0.42, 0.31]
\end{lstlisting}

\subsection{EM fitting}

Po príprave 5D dát sa zavolá funkcia \texttt{\_fit\_em\_standard()},
ktorá GMM nafituje a~priamo vyrenderuje výsledný obraz:

\begin{lstlisting}
def _fit_em_standard(img_np, config, n_comp):
    data_5d, h, w = _prepare_data(img_np)
    gmm = GaussianMixture(
        n_components=n_comp,
        covariance_type='full',
        max_iter=config.get('max_iter', 100),
        init_params='kmeans', random_state=42)
    gmm.fit(data_5d)

    # Transformacia do pixeloveho priestoru
    cov_scale = np.array([[[w**2, w*h], [w*h, h**2]]])
    render_np = render_gaussians(
        gmm.means_[:, :2] * np.array([[w, h]]),  # stredy
        gmm.covariances_[:, :2, :2] * cov_scale, # kovariancie
        np.clip(gmm.means_[:, 2:5], 0, 1),       # farby
        gmm.weights_, (h, w))                     # vahy
    return render_np, gmm, data_5d, h, w
\end{lstlisting}

\noindent
Kľúčové body:
\begin{itemize}[nosep]
  \item Stredy sa transformujú z~normalizovaných súradníc do pixelov
        vynásobením šírkou a~výškou obrazu.
  \item Priestorová časť kovariančných matíc (prvé 2 riadky/stĺpce
        z~$5 \times 5$ matice) sa transformuje do pixelových jednotiek.
  \item Farby sú priamo zo stredov GMM (dimenzie 3--5), orezané do $[0, 1]$.
\end{itemize}

\textbf{Mini-batch variant:} Pre rýchlejšie fitovanie existuje aj
varianta \texttt{\_fit\_em\_minibatch()}, ktorá pracuje len na~náhodnom
podvzorku pixelov (typicky 15\,\%), čím sa výrazne zníži výpočtový čas
za cenu mierne nižšej kvality.

\subsection{Výsledný model}

Výsledkom je GMM, ktorý pre každý pixel vie určiť, s~akou pravdepodobnosťou
patrí do ktorého Gaussiánu.  Každý Gaussián je charakterizovaný:
\begin{itemize}[nosep]
  \item stredom v~priestore $[x, y, r, g, b]$ --- \texttt{gmm.means\_}: $(K, 5)$,
  \item kovariančnou maticou (tvar a~orientácia) --- \texttt{gmm.covariances\_}: $(K, 5, 5)$,
  \item váhou (aký podiel pixelov reprezentuje) --- \texttt{gmm.weights\_}: $(K,)$.
\end{itemize}

Farba Gaussiánu je určená ako strednú hodnotu v~dimenziách~3--5:
$\mathbf{c}_k = \texttt{gmm.means\_}[k,\; 2{:}5]$.

\textbf{Váha Gaussiánu} $\pi_k$ udáva, aký podiel všetkých pixelov v~obraze
je modelovaný daným Gaussiánom.  Súčet všetkých váh je~1.

\subsection{Premenné a~ich rozmery}

\begin{table}[h]
\centering
\begin{tabular}{@{}lll@{}}
\toprule
\textbf{Premenná} & \textbf{Rozmery} & \textbf{Popis} \\
\midrule
\texttt{gt\_np}          & $(H, W, 3)$    & Ground truth obraz, hodnoty v~$[0,1]$ \\
\texttt{render\_np}      & $(H, W, 3)$    & Aproximovaný obraz \\
\texttt{data\_5d}        & $(H{\cdot}W, 5)$ & 5D pixelové dáta \\
\texttt{gmm.means\_}     & $(K, 5)$       & Stredy Gaussiánov \\
\texttt{gmm.covariances\_} & $(K, 5, 5)$  & Kovariančné matice \\
\texttt{gmm.weights\_}   & $(K,)$         & Váhy Gaussiánov \\
\bottomrule
\end{tabular}
\caption{Hlavné premenné a~ich rozmery.}
\label{tab:vars}
\end{table}

% ==========================================================================
\section{Renderovanie obrazu pomocou Gaussiánov}
\label{sec:rendering}
% ==========================================================================

Funkcia \texttt{render\_gaussians} vykresľuje výsledný obraz tak, že
každý Gaussián je \uv{nastriekaný} do pixelového priestoru ako rozmazaná
farebná škvrna.  Implementácia priamo zodpovedá rovnici~\eqref{eq:render}
z~predchádzajúcej sekcie.

\subsection{Kľúčové kroky renderovania}

Nasledujúci kód ukazuje jadro renderovacej slučky:

\begin{lstlisting}
def render_gaussians(means, covs, colors, weights, image_size):
    h, w = image_size
    canvas     = np.zeros((h, w, 3), dtype=np.float32)
    weight_sum = np.zeros((h, w, 1), dtype=np.float32)

    # Mriežka suradnic vsetkych pixelov
    X, Y = np.meshgrid(np.arange(w), np.arange(h))
    pos = np.dstack((X, Y))  # (H, W, 2)

    for mean, cov, color, wgt in zip(means, covs, colors, weights):
        # 1. Vyhodnotenie Gaussianu vo vsetkych pixeloch
        rv  = multivariate_normal(mean=mean, cov=cov)
        pdf = rv.pdf(pos)  # (H, W) -- hustota v kazdom pixeli

        # 2. Peak-normalizacia: maximum = 1 (v strede)
        pdf = pdf / (np.max(pdf) + 1e-8)

        # 3. Vahou skalovany prispevok do platna
        pdf = pdf * wgt
        pdf = np.expand_dims(pdf, axis=-1)  # (H, W, 1)

        canvas     += pdf * color   # farebny prispevok
        weight_sum += pdf           # akumulacia vah

    # Vazeny priemer a orezanie do [0, 1]
    render = np.divide(canvas, weight_sum,
                       out=np.zeros_like(canvas),
                       where=weight_sum > 1e-8)
    return np.clip(render, 0.0, 1.0)
\end{lstlisting}

\noindent
Popis jednotlivých krokov:
\begin{enumerate}[nosep]
  \item \textbf{Vyhodnotenie PDF:} Pre~každý Gaussián sa spočíta hodnota
        normálneho rozdelenia vo~všetkých pixeloch naraz (vďaka mriežke
        \texttt{pos}).
  \item \textbf{Peak-normalizácia:} PDF sa vydelí svojím maximom, aby
        hodnota v~strede Gaussiánu bola~1.  Tým sa oddelí \emph{tvar}
        (kde Gaussián pôsobí) od \emph{intenzity} (váha $\pi_k$).
  \item \textbf{Akumulácia:} Každý Gaussián prispeje do plátna
        farbou $\mathbf{c}_k$ škálovanou váhou a~priestorovým profilom.
  \item \textbf{Finalizácia:} Výsledný obraz je vážený priemer
        príspevkov, orezaný do~$[0, 1]$.
\end{enumerate}

% ==========================================================================
\section{Vyhodnotenie kvality}
\label{sec:evaluation}
% ==========================================================================

Aproximovaný obraz sa porovná s~originálom pomocou týchto metrík:
\begin{itemize}[nosep]
  \item \textbf{MSE} --- Mean Squared Error
  \item \textbf{RMSE} --- Root Mean Squared Error
  \item \textbf{PSNR} --- Peak Signal-to-Noise Ratio (dB)
  \item \textbf{SSIM} --- Structural Similarity Index
  \item \textbf{dSSIM} --- $1 - \text{SSIM}$ (dissimilarity)
  \item \textbf{LPIPS} --- Learned Perceptual Image Patch Similarity (voliteľne)
\end{itemize}

Výsledky sa ukladajú do zoznamu výsledkov a~renderovaný obraz sa uloží ako PNG.

\end{document}
